\begin{thebibliography}{99}
\small

\bibitem{boiko2023} D. A. Boiko, R. MacKnight, B. Kline, and G. Gomes. Autonomous chemical research with large language models. \emph{Nature} 624:570--578, 2023. doi:10.1038/s41586-023-06792-0.

\bibitem{bran2024} A. M. Bran, S. Cox, O. Schilter, C. Baldassari, A. D. White, and P. Schwaller. Augmenting large language models with chemistry tools. \emph{Nature Machine Intelligence} 6:525--535, 2024. doi:10.1038/s42256-024-00832-8.

\bibitem{lu2024} C. Lu et al. The AI Scientist: Towards Fully Automated Open-Ended Scientific Discovery. arXiv:2408.06292, 2024.

\bibitem{gottweis2025} J. Gottweis et al. Towards an AI co-scientist. arXiv:2502.18864, 2025.

\bibitem{kosmos2025} B. Mitchener et al. Kosmos: An AI Scientist for Autonomous Discovery. arXiv:2511.02824, 2025.

\bibitem{lewis2020} P. Lewis et al. Retrieval-Augmented Generation for Knowledge-Intensive NLP Tasks. In \emph{Advances in Neural Information Processing Systems} 33:9459--9474, 2020.

\bibitem{yao2023} S. Yao et al. ReAct: Synergizing Reasoning and Acting in Language Models. In \emph{International Conference on Learning Representations}, 2023.

\bibitem{schick2023} T. Schick et al. Toolformer: Language Models Can Teach Themselves to Use Tools. In \emph{Advances in Neural Information Processing Systems} 36, 2023.

\bibitem{packer2023} C. Packer et al. MemGPT: Towards LLMs as Operating Systems. arXiv:2310.08560, 2023.

\bibitem{mialon2023} G. Mialon et al. Augmented Language Models: a Survey. arXiv:2302.07842, 2023.

\bibitem{park2023} J. S. Park et al. Generative Agents: Interactive Simulacra of Human Behavior. In \emph{Proceedings of UIST 2023}, 2023. doi:10.1145/3586183.3606763.

\bibitem{shinn2023} N. Shinn et al. Reflexion: Language Agents with Verbal Reinforcement Learning. In \emph{Advances in Neural Information Processing Systems} 36, 2023.

\bibitem{doyle1979} J. Doyle. A truth maintenance system. \emph{Artificial Intelligence} 12(3):231--272, 1979. doi:10.1016/0004-3702(79)90008-0.

\bibitem{dekleer1986} J. de Kleer. An assumption-based TMS. \emph{Artificial Intelligence} 28(2):127--162, 1986. doi:10.1016/0004-3702(86)90080-9.

\bibitem{agm1985} C. E. Alchourr\'on, P. G\"ardenfors, and D. Makinson. On the logic of theory change: Partial meet contraction and revision functions. \emph{Journal of Symbolic Logic} 50(2):510--530, 1985. doi:10.2307/2274239.

\bibitem{gardenfors1988} P. G\"ardenfors. \emph{Knowledge in Flux: Modeling the Dynamics of Epistemic States}. MIT Press, 1988.

\bibitem{dung1995} P. M. Dung. On the acceptability of arguments and its fundamental role in nonmonotonic reasoning, logic programming and n-person games. \emph{Artificial Intelligence} 77(2):321--357, 1995. doi:10.1016/0004-3702(94)00041-X.

\bibitem{belnap1977} N. D. Belnap Jr. A Useful Four-Valued Logic. In J. M. Dunn and G. Epstein, editors, \emph{Modern Uses of Multiple-Valued Logic}, pages 5--37. Reidel, 1977.

\bibitem{buneman2001} P. Buneman, S. Khanna, and W.-C. Tan. Why and Where: A Characterization of Data Provenance. In \emph{Database Theory---ICDT 2001}, pages 316--330. Springer, 2001. doi:10.1007/3-540-44503-X_20.

\bibitem{green2007} T. J. Green, G. Karvounarakis, and V. Tannen. Provenance semirings. In \emph{Proceedings of PODS 2007}, pages 31--40, 2007. doi:10.1145/1265530.1265535.

\bibitem{prov2013} W3C. \emph{PROV-O: The PROV Ontology}. W3C Recommendation, 30 April 2013.

\bibitem{wilkinson2016} M. D. Wilkinson et al. The FAIR Guiding Principles for scientific data management and stewardship. \emph{Scientific Data} 3:160018, 2016. doi:10.1038/sdata.2016.18.

\bibitem{mccarthy1993} J. McCarthy. Notes on Formalizing Context. In \emph{Proceedings of the 13th International Joint Conference on Artificial Intelligence}, 1993.

\bibitem{giunchiglia1993} F. Giunchiglia. Contextual reasoning. \emph{Epistemologia}, 1993.

\bibitem{abramsky2011} S. Abramsky and A. Brandenburger. The sheaf-theoretic structure of non-locality and contextuality. \emph{New Journal of Physics} 13:113036, 2011. doi:10.1088/1367-2630/13/11/113036.

\bibitem{curry2014} J. Curry. \emph{Sheaves, Cosheaves and Applications}. PhD thesis, University of Pennsylvania, 2014. arXiv:1303.3255.

\bibitem{robinson2017} M. Robinson. Sheaves are the canonical data structure for sensor integration. \emph{Information Fusion} 36:208--224, 2017. doi:10.1016/j.inffus.2016.12.002.

\bibitem{pearl2009} J. Pearl. \emph{Causality: Models, Reasoning, and Inference}. 2nd edition, Cambridge University Press, 2009.

\bibitem{halpern2016} J. Y. Halpern. \emph{Actual Causality}. MIT Press, 2016.

\bibitem{birkhoff1940} G. Birkhoff. \emph{Lattice Theory}. American Mathematical Society Colloquium Publications, 1940.

\bibitem{ganter1999} B. Ganter and R. Wille. \emph{Formal Concept Analysis: Mathematical Foundations}. Springer, 1999.

\bibitem{cousot1977} P. Cousot and R. Cousot. Abstract Interpretation: A Unified Lattice Model for Static Analysis of Programs by Construction or Approximation of Fixpoints. In \emph{Proceedings of POPL 1977}, pages 238--252, 1977. doi:10.1145/512950.512973.

\bibitem{tarski1955} A. Tarski. A lattice-theoretical fixpoint theorem and its applications. \emph{Pacific Journal of Mathematics} 5(2):285--309, 1955.

\bibitem{hayesroth1985} B. Hayes-Roth. A blackboard architecture for control. \emph{Artificial Intelligence} 26:251--321, 1985. doi:10.1016/0004-3702(85)90063-3.

\bibitem{nii1986} H. P. Nii. The Blackboard Model of Problem Solving and the Evolution of Blackboard Architectures. \emph{AI Magazine} 7(2):38--53, 1986.

\bibitem{baars1988} B. J. Baars. \emph{A Cognitive Theory of Consciousness}. Cambridge University Press, 1988.

\bibitem{dehaene2001} S. Dehaene and L. Naccache. Towards a cognitive neuroscience of consciousness: basic evidence and a workspace framework. \emph{Cognition} 79:1--37, 2001. PMID:11164022.

\bibitem{mashour2020} G. A. Mashour, P. Roelfsema, J.-P. Changeux, and S. Dehaene. Conscious Processing and the Global Neuronal Workspace Hypothesis. \emph{Neuron} 105:776--798, 2020. doi:10.1016/j.neuron.2020.01.026.

\bibitem{bengio2017} Y. Bengio. The Consciousness Prior. arXiv:1709.08568, 2017.

\bibitem{goyal2021} A. Goyal et al. Coordination Among Neural Modules Through a Shared Global Workspace. arXiv:2103.01197, 2021.

\bibitem{gurnee2026} W. Gurnee et al. Verbalizable Representations Form a Global Workspace in Language Models. \emph{Transformer Circuits Thread}, published 6 July 2026; arXiv:2607.15495, posted 16 July 2026.

\bibitem{furnas1987} G. W. Furnas, T. K. Landauer, L. M. Gomez, and S. T. Dumais. The Vocabulary Problem in Human-System Communication. \emph{Communications of the ACM} 30(11):964--971, 1987. doi:10.1145/32206.32212.

\bibitem{swanson1986} D. R. Swanson. Fish oil, Raynaud's syndrome, and undiscovered public knowledge. \emph{Perspectives in Biology and Medicine} 30:7--18, 1986. doi:10.1353/pbm.1986.0087.

\bibitem{garikaparthi2025} A. Garikaparthi, M. Patwardhan, A. S. Kanade, A. Hassan, L. Vig, and A. Cohan. MIR: Methodology Inspiration Retrieval for Scientific Research Problems. In \emph{Proceedings of ACL 2025}, pages 28614--28659. doi:10.18653/v1/2025.acl-long.1390.

\bibitem{carbonell1998} J. Carbonell and J. Goldstein. The Use of MMR, Diversity-Based Reranking for Reordering Documents and Producing Summaries. In \emph{Proceedings of SIGIR 1998}, pages 335--336, 1998. doi:10.1145/290941.291025.

\bibitem{jafarzadeh2021} M. Jafarzadeh et al. A Review of Open World Learning and Steps Toward Open World Learning Without Labels. arXiv:2011.12906, 2021.

\bibitem{saunders2018} B. Saunders et al. Saturation in qualitative research: exploring its conceptualization and operationalization. \emph{Quality \& Quantity} 52:1893--1907, 2018. doi:10.1007/s11135-017-0574-8.

\bibitem{tate2009} R. Tate, M. Stepp, Z. Tatlock, and S. Lerner. Equality Saturation: A New Approach to Optimization. In \emph{Proceedings of POPL 2009}, pages 264--276, 2009. doi:10.1145/1480881.1480915.

\bibitem{willsey2021} M. Willsey et al. egg: Fast and Extensible Equality Saturation. \emph{Proceedings of the ACM on Programming Languages} 5(POPL), 2021. doi:10.1145/3434304.

\bibitem{vanemden1976} M. H. van Emden and R. A. Kowalski. The Semantics of Predicate Logic as a Programming Language. \emph{Journal of the ACM} 23(4):733--742, 1976. doi:10.1145/321978.321991.

\bibitem{rissanen1978} J. Rissanen. Modeling by shortest data description. \emph{Automatica} 14(5):465--471, 1978. doi:10.1016/0005-1098(78)90005-5.

\bibitem{tishby2000} N. Tishby, F. C. Pereira, and W. Bialek. The information bottleneck method. arXiv:physics/0004057, 2000.

\bibitem{fletcher2026} A. H. A. Fletcher and M. Stevenson. Decision-Theoretic Stopping Rules for Document Screening. arXiv:2606.07071, 2026.

\bibitem{fletcher2026b} A. H. A. Fletcher and M. Stevenson. Confidence-Based Stopping Methods for Systematic Reviews. arXiv:2606.15380, 2026.

\bibitem{macfarlane2022} A. MacFarlane, T. Russell-Rose, and F. Shokraneh. Search strategy formulation for systematic reviews: Issues, challenges and opportunities. \emph{Intelligent Systems with Applications}, 2022. doi:10.1016/j.iswa.2022.200091.

\bibitem{shi2017} B. Shi and T. Weninger. Open-World Knowledge Graph Completion. arXiv:1711.03438, 2017.

\bibitem{yang2022} H. Yang, Z. Lin, and M. Zhang. Rethinking Knowledge Graph Evaluation Under the Open-World Assumption. arXiv:2209.08858, 2022.

\bibitem{li2024kg} J. Li, R. Luo, J. Sun, J. Xiao, and Y. Yang. Progressive Knowledge Graph Completion. arXiv:2404.09897, 2024.

\bibitem{gradel2024} E. Gr\"adel and V. Tannen. Provenance Analysis and Semiring Semantics for First-Order Logic. arXiv:2412.07986, 2024.

\bibitem{landau1976} L. D. Landau and E. M. Lifshitz. \emph{Mechanics}. 3rd edition, Pergamon Press, 1976.

\bibitem{strogatz2015} S. H. Strogatz. \emph{Nonlinear Dynamics and Chaos}. 2nd edition, Westview Press, 2015.

\bibitem{dlmf2026} NIST Digital Library of Mathematical Functions. Chapter 19: Elliptic Integrals. National Institute of Standards and Technology, accessed 2026.

\bibitem{reinberger2021} P. Reinberger et al. On the trajectory of the nonlinear pendulum: Exact analytic solutions via power series. arXiv:2108.09395, 2021.

\bibitem{graber2018} J. Graber-Mitchell. Finding the period of a simple pendulum. arXiv:1805.00002, 2018.

\bibitem{nosek2015} Open Science Collaboration. Estimating the reproducibility of psychological science. \emph{Science} 349(6251):aac4716, 2015. doi:10.1126/science.aac4716.

\bibitem{munafo2017} M. R. Munaf\`o et al. A manifesto for reproducible science. \emph{Nature Human Behaviour} 1:0021, 2017. doi:10.1038/s41562-016-0021.

\end{thebibliography}
